\section{관련 연구}
\label{sec:related-work}

\subsection{트랜스포머}
\label{sec:rw-transformer}

시퀀스 모델링(sequence modelling)의 지형은 트랜스포머 아키텍처~\citep{vaswani2017attention}의 등장으로 근본적으로 다시 짜였다 — 트랜스포머는 순환 레이어(recurrent layer)를 버리는 대신 병렬화 가능한 셀프 어텐션(self-attention) 메커니즘을 도입했다.
이후 분야는 BERT~\citep{devlin2019bert}처럼 판별적 과제에 최적화된 인코더 전용 모델과, GPT-3~\citep{brown2020language}처럼 거대한 스케일링(scaling)과 창발적인 인-컨텍스트 학습(in-context learning)으로 현재의 생성형 AI 시대를 견인한 디코더 전용(decoder-only) 아키텍처로 갈라졌다.
이어진 연구는 시각 인식을 위한 Vision Transformer(ViT)~\citep{dosovitskiy2021image}와, 텍스트-투-텍스트 처리의 통일을 시도한 T5 프레임워크~\citep{raffel2020exploring} 등으로 아키텍처의 적용 범위를 넓혀 왔다.
최근의 진전은 계산 효율성과 멀티모달리티(multimodality)에 무게를 두며, 특히 FlashAttention~\citep{dao2022flashattention} 같은 하드웨어 인지 최적화와, Mixtral $8{\times}7$B~\citep{jiang2024mixtral}처럼 Mixture-of-Experts(MoE)~\citep{fedus2022switch}를 채택한 모델에서 두드러진다.
현재 패러다임에서는 Gemini~1.5~\citep{gemini2024gemini}와 GPT-4o~\citep{hurst2024gpt} 같은 모델이 합성형(compositional) 아키텍처를 넘어 네이티브 멀티모달리티(native multimodality)로 옮겨 가, 이질적인 데이터 스트림을 가로지르는 매끄러운 추론을 가능하게 했다.

이 지형 안에서 PRAGMA는 \emph{이질적 표 형식 이벤트 스트림을 위한 인코더형 파운데이션 모델}로 이해될 수 있다. 금융 거래에서 출발한 동기이지만, 엔티티가 시간을 따라 불규칙하고 다중 필드 레코드를 누적하는 어떤 도메인에도 자연스럽게 확장된다.
PRAGMA는 인코더 전용 트랜스포머의 확장성과 양방향 맥락화를 그대로 물려받으면서, 이를 이질적 필드, 명시적 시간 신호, 재사용 가능한 레코드 수준 표현에 맞춰 변형해 적용한다.

\subsection{마스킹 모델링}
\label{sec:rw-masked-modelling}

생성형 디코더의 스케일링과 나란히, 마스킹 모델링은 자기지도 표현 학습(self-supervised representation learning)의 지배적 패러다임으로 자리잡았다.
이는 양방향 맥락을 포착하는 마스킹 언어 모델링(masked language modelling, MLM) 목적함수를 도입한 BERT~\citep{devlin2019bert}로 시작되었으며, RoBERTa~\citep{liu2019roberta}는 동적 마스킹(dynamic masking)과 최적화된 학습 레시피(recipe)로 이를 한층 다듬었다.
이후 MLM의 성공은 시각 도메인의 마스킹 이미지 모델링(masked image modelling, MIM)으로 옮겨졌고, BEiT~\citep{bao2021beit}와 Masked Autoencoders(MAE)~\citep{he2022masked}는 가려진 이미지 패치(image patch)를 복원하는 과정이 모델로 하여금 전체 구조적 표현을 학습하게 만든다는 것을 보였다.
최근에는 Data2Vec~\citep{baevski2022data2vec}처럼 모달 간 통합(cross-modal unification)으로, 그리고 Joint-Embedding Predictive Architecture(I-JEPA)~\citep{assran2023self}처럼 원시 신호 복원에서 잠재 피처 예측(latent feature prediction)으로 이동하는 흐름이 나타나고 있다.

PRAGMA는 이 흐름에서 직접 영감을 받지만, 마스킹 모델링을 텍스트와 이미지에서 \emph{이질적 금융 레코드}로 확장한다.
저자들의 목적함수는 개별 토큰, 이벤트 전체, 의미 타입(키)을 마스킹하며, 부분 관측 이벤트의 복원과 전체 거래 이력으로부터의 전이 가능한 표현 학습을 함께 유도한다.



\subsection{표 형식 데이터를 위한 트랜스포머}
\label{sec:rw-tabular}

구조화 데이터(structured data)에서는 그동안 그래디언트 부스팅 결정 트리(GBDT)가 지배적이었지만, 트랜스포머는 새로운 부류의 ``표 형식 딥러닝(Tabular Deep Learning)'' 아키텍처를 촉발했다.
TabTransformer~\citep{huang2020tabtransformer}와 FT-Transformer~\citep{gorishniy2021revisiting} 같은 초기 모델은 셀프 어텐션으로 피처 간 의존을 모델링하는 데 집중해, 고차원 데이터셋에서 GBDT와 동등한 성능을 보였다.
SAINT~\citep{somepalli2021saint}는 피처와 행 상호작용을 모두 다루는 이중 어텐션을 도입했고, Trompt~\citep{chen2023trompt}는 표의 본질 속성과 표본별 변동을 분리하기 위한 프롬프트 튜닝을 제안했다.
TabPFN~\citep{hollmann2023tabpfn}은 베이즈 추론을 근사하도록 합성 데이터로 사전학습된 파운데이션 모델로 패러다임 전환을 가져왔다.
인-컨텍스트 학습을 활용해, TabPFN은 단일 순전파만으로 예측을 만들어 내며 반복 학습의 필요를 제거한다.
원래 모델은 1{,}000개 표본으로 제한되었지만, TabPFN-v2와 TabPFN-v2.5~\citep{hollmann2025accurate,grinsztajn2025tabpfn}는 100{,}000개 표본과 실제 데이터의 복잡성을 다루도록 아키텍처를 확장하고, 범주형 피처·결측값·이상치를 기본 지원한다.
가장 최근의 Mitra~\citep{zhang2025mitra}는 SAINT의 이중 어텐션을 채택하면서도, 거대한 합성 사전(prior) 혼합으로만 사전학습되는 TabPFN의 파운데이션 모델 패러다임을 따른다.

PRAGMA는 필드 정체성을 보존하고 어텐션으로 필드 간 상호작용을 모델링한다는 점에서 표 형식 트랜스포머와 정신적으로 가깝지만, TabTransformer·FT-Transformer·SAINT와 달리 고정 스키마의 단일 행 위에서 작동하지 않는다.
합성된 지도(supervised) 과제로 학습되는 TabPFN류 표 형식 파운데이션 모델과 비교하면, PRAGMA는 실제 금융 원장 위에서 자기지도로 사전학습되며 계층형 인코더로 이질적 이벤트의 가변 길이 사용자 이력을 모델링한다.

\subsection{추천 시스템 모델링}
\label{sec:rw-recsys}

순차 추천 모델은 풍부한 부가 정보를 동반한 시간순 이벤트 시퀀스를 다룬다는 점에서, 거래 모델링과 구조적 유사성을 갖는다.
트랜스포머 기반 추천 모델은 사용자 상호작용 이력을 토큰 시퀀스로 다룬다 — SASRec~\citep{kang2018self}는 순환 대신 셀프 어텐션으로 장거리 의존(long-range dependency)을 포착했고, BERT4Rec~\citep{sun2019bert4rec}는 마스킹된 아이템 예측을 통한 양방향 맥락화(bidirectional contextualisation)가 더 견고한 표현을 만들어 냄을 보였다.
이후 분야는 LLM 패러다임으로 수렴했다 — P5~\citep{geng2022recommendation}는 다양한 추천 과제를 T5 기반의 통일된 텍스트-투-텍스트(text-to-text) 프레임워크로 환원했고, TALLRec~\citep{bao2023tallrec}은 명령 튜닝(instruction tuning)을 도입해 범용 LLM을 추천 로직에 정렬(align)했다.

더 최근의 산업적 흐름은, 양성 상호작용만 모델링하는 데에서 더 풍부한 이벤트 스트림을 인코딩하는 쪽으로 이동했다.
Generative Recommenders~\citep{zhai2024actions}는 인과 시퀀스(causal sequence) 안에 아이템 토큰과 액션 토큰을 교차 배치해, 멱법칙(power-law)에 가까운 품질 향상으로 수조 파라미터 규모로 확장한다.
ARGUS~\citep{khrylchenko2025scaling}는 자기회귀(autoregressive) 학습을 피드백 예측과 다음 아이템 예측으로 분해하여, 추천 트랜스포머를 10억 파라미터 규모로 키운다.
TransAct 계열 연구~\citep{xia2023transact,xia2025transact}는 각 사용자 행동을 콘텐츠·행동 유형·맥락의 조합으로 임베딩해 CTR(click-through rate) 예측을 수행하고, 평생(lifelong) 행동 시퀀스로 확장한다.

PRAGMA는 시간순 이벤트 이력과 자기지도 사전학습을 사용한다는 점에서 이 문헌과 가깝다.
다만 각 상호작용을 흔히 아이템 토큰으로 환원하는 추천 모델과 달리, PRAGMA는 타입이 있는 필드, 금액, 자유 텍스트, 시간 좌표를 갖춘 더 풍부한 금융 이벤트를 모델링하며, 랭킹을 넘어선 폭넓은 뱅킹 과제에 적응된다.

\subsection{금융을 위한 파운데이션 모델}
\label{sec:rw-finance}

금융 파운데이션 모델 패러다임은 특화된 텍스트 인코더에서 다양한 데이터 모달리티를 통합하는 종합 추론 엔진으로 빠르게 성숙해 왔다.
이런 진화는 인코더 전용 아키텍처를 금융 코퍼스에 적응시킨 FinBERT~\citep{yang2020finbert}로 시작되어, 감성 분석(sentiment analysis)이나 ESG 분류 같은 판별적 과제의 엄격한 베이스라인을 제시했다.
이후 분야는 BloombergGPT~\citep{wu2023bloomberggpt}와 함께 거대한 생성형 스케일로 옮겨 갔는데, 이는 독점 금융 데이터셋과 일반 웹 코퍼스를 교차해 학습하면 도메인 특화 벤치마크에서 우월한 성능을 낸다는 점을 보였다.
이런 거대 모델의 접근성 장벽을 다루기 위해, FinGPT~\citep{yang2023fingpt}는 데이터 중심의 가벼운 적응 프레임워크를 제안하며, 오픈소스 모델의 효율적 LoRA 파인튜닝~\citep{hu2022lora}을 통해 금융 LLM의 접근성을 민주화했다.
가장 최근에는 텍스트의 경계를 넘어 시장 데이터의 구조적 성격을 다루는 연구가 나타났으며, Time-LLM~\citep{jin2024timellm}과 Chronos~\citep{ansari2024chronos} 같은 모델은 수치 시계열(numerical time series)을 토큰 시퀀스로 다루어, 트랜스포머가 제로샷 예측(zero-shot forecasting)을 수행할 수 있게 한다.

이런 구조적 전환을 소비자 금융으로 확장하면서, 최근의 파운데이션 모델은 대규모 사용자 거래 원장 위에서 직접 학습되고 있다.
예컨대, nuFormer~\citep{braithwaite2025your}는 토큰화된 거래 시퀀스를 전통적인 표 형식 피처와 결합하면 실제 위험 예측에서 수작업 피처 엔지니어링을 효과적으로 대체할 수 있음을 보인다.
같은 흐름에서 TransactionGPT~\citep{dou2025transactiongpt}는 십억 단위 결제 궤적의 멀티모달, 시간, 표 형식 차원을 명시적으로 모델링하기 위한 특화 3D-Transformer 아키텍처를 도입해, 다운스트림 이상 탐지와 궤적 생성에서 최첨단 성능을 달성한다.

PRAGMA는 주로 금융 언어를 다루는 FinBERT, BloombergGPT, FinGPT 같은 텍스트 중심 금융 파운데이션 모델, 그리고 수치 시계열을 토큰화해 예측에 쓰는 Time-LLM과 Chronos와는 다르다.
nuFormer와 TransactionGPT 같은 거래 원장 모델에 더 가깝지만, PRAGMA는 명시적 프로필 상태가 있는 멀티 소스 뱅킹 이벤트 위의 \emph{재사용 가능한 인코더 백본}을 지향하며, 다양한 판별적 과제에 대해 가벼운 적응을 제공한다.
